\chapter{Especificaciones Técnicas de Montaje e Instalación}

\section{Procedimientos de Montaje de Bandejas Portacables}
1. El trazado de la ruta de bandejas se coordinará previamente con las estructuras metálicas de la nave para evitar interferencias con las cerchas y puente grúa.
2. Los soportes se fijarán mediante tacos expansivos de acero en vigas de concreto o fijación soldada en estructuras metálicas, espaciados a no más de $2.00\text{ m}$.
3. Todas las secciones de bandeja se unirán con eclisas de conexión y pernos de cabeza redonda con tuerca de freno. Se asegurará la continuidad eléctrica instalando puentes de masa en cable de $6\text{ mm}^2$ verde/amarillo entre tramos adyacentes.

\section{Tendido e Instalación de Conductores}
1. Antes de proceder al cableado, se comprobará la limpieza interna y secado de todas las tuberías mediante soplado o guía de nylon.
2. Se utilizarán lubricantes sintéticos neutros especiales para cables libres de halógenos.
3. El esfuerzo de tracción al jalar los cables no superará los $5.0\text{ kg/mm}^2$ de la sección total de cobre.
4. Se dejarán bucles de reserva de $1.00\text{ m}$ en los tableros generales y $0.30\text{ m}$ en las cajas de paso o derivados.

\section{Instalación y Pruebas del Sistema de Puesta a Tierra (SPAT)}
1. Se excavará la zanja perimetral de $0.40\text{ m}$ de ancho por $0.80\text{ m}$ de profundidad en el perímetro de la nave.
2. Se colocarán las 6 varillas Copperweld de $5/8'' \times 2.40\text{ m}$ hincadas verticalmente.
3. Se tenderá el cable de cobre desnudo de $35\text{ mm}^2$ interconectando todas las varillas y el tablero TGD.
4. Se rellenará la zanja aplicando compuesto mejorador de conductividad gel (Thorjel / Bentonita) en capas compactadas de $0.20\text{ m}$.
5. Se medirá la resistencia final mediante el **método de caída de potencial de las 3 varillas (NTP 370.304)**, certificando un valor menor a $5.0 \:\Omega$.

\section{Pruebas de Aislamiento y Puesta en Servicio}
1. **Prueba de Resistencia de Aislamiento:** Medición con megóhmetro a $1000\text{V DC}$ entre fases, fase-neutro y fase-tierra. El valor mínimo aceptable será de $50\text{ M}\Omega$.
2. **Prueba de Continuidad:** Verificación de baja resistencia en todos los conductores de protección (PE) desde los tableros hasta las carcasas metálicas de los motores y estructuras.
3. **Prueba de Funcionamiento de Protecciones Diferenciales:** Verificación del tiempo de disparo ($\le 30\text{ ms}$) mediante simulador de corriente de defecto diferencial.

\chapter{Cronograma de Ejecución de Obra}

\section{Plan de Trabajo}
La ejecución de las obras de las instalaciones eléctricas industriales está programada para un período de **45 días calendarios**, distribuidos en las siguientes etapas consecutivas:

\begin{table}[H]
\centering
\caption{Cronograma de ejecución de partidas de la nave industrial}
\begin{tabular}{clc}
\toprule
\textbf{Semana} & \textbf{Actividad / Partida Principal} & \textbf{Avance \%} \\
\midrule
Semana 1 & Movilización, trazo, replanteo y apertura de zanjas SPAT & $10\%$ \\
Semana 2 & Hincado de varillas, tendido de malla de cobre e instalación de SPAT & $25\%$ \\
Semana 3 & Montaje de bandeja portacables principal a $7.00\text{ m}$ de altura & $45\%$ \\
Semana 4 & Instalación de tuberías empotradas PVC SAP y tableros TGD/TF1/TI1 & $65\%$ \\
Semana 5 & Cableado de energía N2XH, montaje de High-Bay LED y tomas Stecker & $85\%$ \\
Semana 6 & Conexión de motores, banco de condensadores, pruebas e inspección & $100\%$ \\
\bottomrule
\end{tabular}
\end{table}
